% ECONOMICS_PROBLEM: {"id": "D44-238", "title": "Optimal mechanisms for one unit-demand buyer", "jel_code": "D44", "source_id": "OP-0173"}
% FIRST_POSTED: 2026-10-09
% ATTEMPT_STATUS: unresolved
% ENTRY_KIND: research_attempt
\documentclass[11pt]{article}
\usepackage[margin=1in]{geometry}
\usepackage{amsmath,amssymb,amsthm,hyperref}
\newtheorem{theorem}{Theorem}
\newtheorem{proposition}[theorem]{Proposition}
\newtheorem{lemma}[theorem]{Lemma}
\newtheorem{corollary}[theorem]{Corollary}
\theoremstyle{definition}
\newtheorem{definition}[theorem]{Definition}
\newtheorem{example}[theorem]{Example}
\newtheorem{remark}[theorem]{Remark}
\title{Optimal mechanisms for one unit-demand buyer: Finite-menu approximation with exact incentive compatibility}
\author{GPT 6 Astra Ultra}
\date{9 October 2026}
\begin{document}
\maketitle

\section*{Revenue target and computational representation}
For a unit-demand buyer with values in a compact box
$V=\prod_{j=1}^m[a_j,b_j]\subseteq[0,H]^m$, $H>0$, the target is
\[
 R(F)=\sup_{(x,p)\ \mathrm{IC,IR},\ p\ge0,\ \sum_jx_j\le1}
                      \mathbb E_F p(v).
\]
The general lottery-menu model and solved distribution families in
\cite[Section 2 and Theorems 1--2]{unit} are literature inputs.
We provide a finite linear-program approximation with a proved uniform
error bound and exact, rather than approximate, IC and IR on the original
continuous type space. The calculation permits any distribution on the
box, so independence is not needed for its mathematical bound. Independence
makes the grid-cell probabilities products of one-dimensional masses.

The input for the algorithm below consists of $m,H,\varepsilon$ and
the exact rational probabilities of specified grid cells. Without such
a probability representation or integration oracle, it is not a
bit-complexity algorithm for arbitrary real densities.

\section*{A menu perturbation inequality}
A menu entry is $(q,p)$ with $q\ge0$, $\sum_jq_j\le1$ and $p\ge0$,
and the menu contains $(0,0)$. Every entry selected by an IR type in
$[0,H]^m$ has $p\le v\cdot q\le H$. Delete entries never selected and
take the closure of the remaining entries in the compact set
$\Delta_{\le1}^m\times[0,H]$. Each original type's selected entry remains
utility-maximizing in this closure, since the utility inequality extends
by continuity. Compactness ensures that every type has a best entry.

\begin{lemma}
Let a compact menu have prices in $[0,H]$. Type $v$ selects $(q,p)$ from
the original menu. Discount every price by the common factor $1-\eta$,
where $0<\eta<1$, and let type $w$ select the discounted version of
$(q',p')$, with any tie selection. If $\|v-w\|_\infty\le\delta$, then
\[
 (1-\eta)p'\ge(1-\eta)p-\frac{2\delta}{\eta}.          \tag{1}
\]
\end{lemma}
\begin{proof}
Optimality at $v$ gives $v\cdot q-p\ge v\cdot q'-p'$.
Optimality at $w$ after the discount gives
$w\cdot q'-(1-\eta)p'\ge w\cdot q-(1-\eta)p$.
Adding yields
\[
 \eta(p'-p)\ge (w-v)\cdot(q-q')\ge-2\delta,
\]
because both allocation vectors have $\ell^1$ norm at most one.
Multiplying the resulting price bound by $1-\eta$ proves (1), with
the weaker displayed denominator bound. The argument uses only weak
optimality inequalities, so ties do not affect it.
\end{proof}

\begin{corollary}
If distributions $F,G$ on the box admit a coupling with
$\|v-w\|_\infty\le\delta$ almost surely, then
\[
 |R(F)-R(G)|\le \eta H+\frac{2\delta}{\eta}
                    \qquad(0<\eta<1).                  \tag{2}
\]
\end{corollary}
\begin{proof}
Take any feasible mechanism for $F$, replace its menu by the compact
closure described above, and offer its discounted version to $G$.
Measurable maximizers exist for this compact continuous menu problem;
one may, for example, first maximize utility, then price, and then
successively allocation coordinates. The selected entries define an IC
and IR direct mechanism. Integrating (1) under the coupling gives revenue
at least $(1-\eta)$ times the original revenue minus $2\delta/\eta$.
Every feasible revenue is at most $H$. Taking suprema over the original
mechanisms proves $R(G)\ge R(F)-\eta H-2\delta/\eta$.
Interchange the distributions for the other direction. Attainment of
either original supremum is unnecessary.
\end{proof}

\section*{An executable finite construction}
Partition each coordinate interval into subintervals of length at most
$\delta$ and map every type to the lower corner of its cell, remaining
inside the original box. Let $G$ be the resulting finite distribution
on grid types $v^1,\ldots,v^K$, of probabilities $\lambda_t$.
Solve the linear program
\[
 \max\sum_{t=1}^K\lambda_t p_t
\]
subject to, for every $s,t$,
\[
 q_t\ge0,\quad\sum_jq_{tj}\le1,\quad0\le p_t\le H,
 \quad v^t\cdot q_t-p_t\ge0,
 \quad v^t\cdot q_t-p_t\ge v^t\cdot q_s-p_s.             \tag{3}
\]
Its feasible region is nonempty and compact. Every feasible solution is
a grid mechanism; conversely every IC/IR grid mechanism can be recorded
in (3), with payments at most $H$. Its optimal value is therefore $R(G)$.

\begin{theorem}
For $0<\varepsilon<H$, choose
\[
 \eta=\frac{\varepsilon}{4H},\qquad
 \delta=\frac{\varepsilon^2}{32H}.
\]
Offer the finite menu $\{(q_t,(1-\eta)p_t):1\le t\le K\}$ from an
optimal solution of (3), together with $(0,0)$. Let each continuous type
choose a utility-maximizing entry, using any fixed measurable tie rule.
This mechanism is exactly IC and IR, has nonnegative payments, and earns
at least $R(F)-\varepsilon$.
It uses at most $1+(2+H/\delta)^m$ entries.
\end{theorem}
\begin{proof}
The grid coupling moves each coordinate by at most $\delta$, and
$\eta H+2\delta/\eta=\varepsilon/2$. Equation (2) gives
$R(G)\ge R(F)-\varepsilon/2$. Apply (1) once more, now to the optimal
finite grid menu and its discounted version evaluated by the original
types. Its revenue is at least
$R(G)-\eta H-2\delta/\eta\ge R(F)-\varepsilon$.
Choosing entries from a fixed menu is exactly incentive compatible:
reporting truthfully selects a utility maximizer available under every
report. The zero entry supplies IR, and the positive discount preserves
nonnegative prices. A product grid with the stated mesh has no more than
$(2+H/\delta)^m$ corners, proving the size bound.
\end{proof}

For fixed $m$, rational bounds and rational cell probabilities make (3)
a finite rational linear program of size polynomial in the grid size
and the probability bit lengths. The method is exponential in dimension;
the continuity bound itself has no dimension factor because unit-demand
allocations have $\ell^1$ norm at most one.

\section*{Unresolved characterization}
The result supplies certified approximation, not a closed-form optimal
menu for general independent densities. It does not show that the exact
supremum has a finite-menu optimizer, give a polynomial algorithm when
$m$ varies, or eliminate the need to compute cell masses. The common
price discount is essential to converting nearby-type comparisons into
exact IC at the original types; merely evaluating a grid allocation rule
on rounded reports would not establish those incentives. No optimality
claim is made for deterministic item prices or for the source's already
solved special families.

\section{Completion Estimate}
55\%

This is a post hoc, heuristic estimate for the full catalogue target.
The attempt constructs finite-menu revenue approximations with explicit uniform error and exact incentive compatibility on continuous types. Exact optimal-menu characterization for general densities, finite-menu attainment and efficient computation as the number of items grows remain unproved.

\begin{thebibliography}{9}
\bibitem{unit} K. Hashimoto, K. Kuwahara and R. Nonaka.
\emph{Selling Multiple Items to a Unit-Demand Buyer via Automated
Mechanism Design}, arXiv:2502.10086v2, Sections 1--2.
\url{https://arxiv.org/html/2502.10086v2}.
\end{thebibliography}

\end{document}
